% TOP_PROBLEM: {"id":30007549,"problem_number":"LOCAL-30007549","title":"Innovation rents and competition","category_id":21,"category":{"id":21,"name":"economics","display_name":"Economics","description":"Open economic theory statements, published research agendas, and proposed scoped specifications; question provenance and historical priority are qualified per record.","slug":"economics","order_index":21,"created_at":"2026-10-08T00:00:00Z"},"status":"research_question","external_url":null,"legacy_ids":[],"published":false,"statement":"Find optimal patent-duration and competition-policy choices in a declared dynamic innovation economy with endogenous entry, research, and diffusion.","economics":{"original_id":38,"field":"Growth and productivity","kind":"design","formalization_basis":"adapted_research_question","source_status":"explicit_open","priority_status":"earliest_found","priority_note":"This is the earliest explicit relevant statement verified in this audit, not a claim of original priority; older literature and earlier versions may contain antecedents.","earliest_verified_reference":{"authors":["Ufuk Akcigit"],"title":"Creative destruction in economic growth","year":2026,"url":"https://onlinelibrary.wiley.com/doi/10.1111/sjoe.70037","doi":"10.1111/sjoe.70037","locator":"Section 7, “Open questions,” item (3)","evidence_summary":"The author explicitly identifies unresolved balancing of competition with innovation rents and asks how policy should support innovation without protecting stagnating incumbents."},"current_reference":{"authors":["Ufuk Akcigit"],"title":"Creative destruction in economic growth","year":2026,"url":"https://onlinelibrary.wiley.com/doi/10.1111/sjoe.70037","doi":"10.1111/sjoe.70037","locator":"Section 7, “Open questions,” item (3)","evidence_summary":"The author explicitly identifies unresolved balancing of competition with innovation rents and asks how policy should support innovation without protecting stagnating incumbents."},"formalization_note":"The source explicitly identifies balancing competition and innovation rents as open. This patent-duration/entry-policy optimization is a scoped adaptation."},"statusQualification":"Published question or research agenda verified at the cited locator. The mathematical specification is adapted for this catalogue and includes additional assumptions; source publication does not establish first-ever priority or certify this exact model remains unsolved. The source explicitly identifies balancing competition and innovation rents as open. This patent-duration/entry-policy optimization is a scoped adaptation.","statusReviewedAt":"2026-10-08"}
% FIRST_POSTED: 2026-10-08
% ENTRY_KIND: statement_only
% !TeX program = lualatex
\documentclass[11pt]{article}
\usepackage{fontspec}
\usepackage[a4paper,margin=25mm]{geometry}
\usepackage{amsmath,amssymb,mathtools}
\usepackage{xurl}
\usepackage[hidelinks]{hyperref}
\newcommand{\sref}[2]{\hyperlink{source:#1:#2}{[S#2]}}
\setlength{\emergencystretch}{3em}
\begin{document}
% problemId:30007549
\section*{Innovation rents and competition}\label{30007549}
\subsection{Definitions and mathematical statement}
Fix a Schumpeterian product-line economy in which incumbents and entrants choose research effort and successful innovations improve quality by a specified amount. Patent protection grants innovation rents for duration \(T\), while competition policy \(a\) determines a declared entry restriction or merger rule. Researchers are scarce, consumers have specified preferences, and diffusion produces an explicitly modeled knowledge spillover. Let \(\mathcal P\) be a compact set of feasible pairs \((T,a)\), excluding unspecified discretionary interventions.
For an initial state and discount factor \(\beta\in(0,1)\), define social welfare
\[
W(T,a)=E_{T,a}\sum_{t\ge0}\beta^t
u(C_t).
\]
Here \(C_t\) is aggregate consumption after research and enforcement resource use, and \(u\) is household utility; these costs enter resource feasibility rather than being subtracted twice. Characterize the maximizing feasible policy and identify the response parameters needed to rank alternatives. Include induced entry, incumbent research, replacement, diffusion, and equilibrium research wages; treating innovation rates as fixed while changing rents misses the central tradeoff. Determine which protections reward socially valuable innovation and which mainly preserve obsolete incumbents. Validate predicted innovation responses using credible policy variation and distinguish rent measurement from technological quality. Report how the optimum changes with spillovers, enforcement, and privately known quality. This is a conditional dynamic policy-design target, not a universal assertion that stronger competition always raises or lowers innovation.

\par\textbf{Formulation and scope.} The source explicitly identifies balancing competition and innovation rents as open. This patent-duration/entry-policy optimization is a scoped adaptation.

\par\textbf{Open-question provenance.} An explicit question or research agenda was verified at the source locator. This does not by itself certify that every later version remains unresolved \sref{30007549}{1}.

\par\textbf{Historical priority.} This is the earliest explicit relevant statement verified in this audit, not a claim of original priority; older literature and earlier versions may contain antecedents.

\subsection{Short English statement}
Find optimal patent-duration and competition-policy choices in a declared dynamic innovation economy with endogenous entry, research, and diffusion.

\subsection{Sources}
\begin{itemize}\raggedright
\item[\textnormal{[S1]}]\hypertarget{source:30007549:1}{} Ufuk Akcigit. 2026. Creative destruction in economic growth. Section 7, “Open questions,” item (3).
\url{https://onlinelibrary.wiley.com/doi/10.1111/sjoe.70037}.
DOI: 10.1111/sjoe.70037.
{\small\textit{Earliest verified question/agenda reference; historical priority qualified below. The author explicitly identifies unresolved balancing of competition with innovation rents and asks how policy should support innovation without protecting stagnating incumbents.}}
\end{itemize}
\end{document}
